\Needspace{6\baselineskip}
\begin{origexamplebox}{例题 9.1【大小圆环】}
两个同心导电圆环放置在同一平面上，回路电阻大小均为 $r$。

(1) 若大圆环中突然通入顺时针电流，求小圆环中感生电流的方向；

(2) 若小圆环中突然通入顺时针电流，求大圆环中感生电流的方向。

% Pure source-question geometry: no derived forces, currents or solution curves.
\begin{center}\begin{tikzpicture}[scale=.68,>=Latex]
 \foreach \x/\case in {0/1,7/2}{\begin{scope}[shift={(\x,0)}]
 \draw[thick](0,0)circle(2);\draw[DeepBlue,thick](0,0)circle(1.1);
 \ifnum\case=1\draw[vector](2,0)arc[start angle=0,end angle=-65,radius=2];\node at(0,2.5){大环电流突然建立};
 \else\draw[vector](1.1,0)arc[start angle=0,end angle=-65,radius=1.1];\node at(0,2.5){小环电流突然建立};\fi
 \node[below]at(0,-2.3){两环电阻均为 $r$};\end{scope}}
\end{tikzpicture}\end{center}
\figcap{两种激励情形；箭头仅表示题设中主动建立的顺时针电流。}\label{fig:em24-question-1}
\end{origexamplebox}

\begin{xesolution}
规定垂直纸面向外为磁通正方向。

(1) 大环的顺时针电流使穿过小环的向里磁通从零增大。小环要用向外的磁场阻碍这个增量，故其感生电流
为逆时针。

(2) 互感系数具有互易性，$M_{12}=M_{21}$。小环的顺时针电流同样使大环中的向里磁通增大，所以大环
中的感生电流也为逆时针。结果与谁是“大环”无关，只由磁通的变化方向决定。

图~\ref{ch24:fig:review-concentricloopsdiagram-1}中，所激励环的顺时针电流增加时，另一环的感应电流为逆时针；交换激励环仍须重算所穿磁通。
\par\noindent\begin{minipage}{\linewidth}
\concentricLoopsDiagram
\figcap{同心大小圆环的感应方向}\label{ch24:fig:review-concentricloopsdiagram-1}
\end{minipage}\par\medskip
\end{xesolution}

\Needspace{6\baselineskip}
\begin{origexamplebox}{例题 9.2【奇怪的圆环】}
一个条形磁铁的周围有三个导体圆环，分别位于 $S$ 极左侧、磁铁中部和 $N$ 极右侧，三环与磁铁同轴。
若磁铁突然向左运动，判断三个圆环的加速度方向，以及它们各自有扩张还是收缩的趋势。

% Pure source-question geometry: no derived forces, currents or solution curves.
%
\begin{center}\begin{tikzpicture}[scale=0.8,>=Latex]
  \draw[DeepBlue,dashed,very thick] (0,1.25) arc(90:270:.33 and 1.25);
  \draw[fill=SoftGray] (-1.8,-0.65) rectangle (1.8,0.65);
  \node at (-1.25,0) {$S$}; \node at (1.25,0) {$N$};
  \foreach \x in {-3.4,3.4}{\draw[DeepBlue,very thick] (\x,0) ellipse (0.33 and 1.25);}
  \draw[DeepBlue,very thick] (0,-1.25) arc(270:450:.33 and 1.25);
  \draw[vector] (-1.7,1.1)--(-2.7,1.1) node[midway,above=4pt] {$\vect v_M$};
  \end{tikzpicture}\end{center}
\figcap{磁铁开始左移，三个导体圆环的运动与形变待求。}\label{fig:em24-question-2}
\end{origexamplebox}

\begin{xesolution}
轴向力与径向张力必须分别判断，不能用环内总磁通的方向代替环线所在处的局部磁场方向。
忽略自感、初始环静止时，设磁铁向左速率为 $V>0$，以磁铁中心为原点、向右为 $z$ 正向，
记半径固定的环所穿正向磁通为 $\Phi(z)$、电阻为 $R_c$，则
\[
 I=-\frac{V}{R_c}\frac{\dd\Phi}{\dd z},\qquad
 F_z=I\frac{\dd\Phi}{\dd z}=-\frac V{R_c}\left(\frac{\dd\Phi}{\dd z}\right)^2\le0.
\]
所以非对称位置的圆环初始均向左加速；若保留自感，电流需要建立，严格初始瞬间可仍为零。

左环面对正在接近的 $S$ 极，穿过它的磁通增大。感应电流使圆环与接近的磁极相斥，因此左环向左加速；其感应磁矩
与当地轴向外磁场反向。在环线处的场仍朝右的通常近轴极外位置，径向 Ampere 力向内，故收缩。

右环面对正在远离的 $N$ 极，原磁通减小。它用同向磁矩维持磁通，与磁铁相吸，所以向左加速；
在环线处轴向场朝右时，径向力向外，故扩张。若环很大，线处的轴向场可能反向，就不能沿用这个径向结论。

中环在完全对称的初始位置上有 $\partial\Phi/\partial x=0$，所以理想化的瞬间感应电流和加速度都可为零。磁铁稍微左移后，
通过中环的磁通减小，中环会向左追随磁铁。在磁铁外侧的近赤道区域，局部轴向场与磁矩相反，
因此中环稍后通常是收缩，不是扩张。

可用磁偶极模型检验这一点。设磁矩大小为 $\mu_M$、环半径为 $a$，则
\[
 \Phi(z)=\frac{\mu_0\mu_Ma^2}{2(a^2+z^2)^{3/2}},\quad
 B_z(a,z)=\frac{\mu_0\mu_M(2z^2-a^2)}{4\pi(a^2+z^2)^{5/2}}.
\]
对 $z>0$ 有 $\Phi'<0$，感应电流为正；在 $0<z<a/\sqrt2$ 时局部 $B_z<0$，
每段径向力 $\dd F_a=I B_z\dd\ell<0$，即收缩；更远处才变为扩张。
这也说明没有环半径和位置等几何条件时，应给出条件分支，而不能把三个径向趋势当作普遍结论。

图~\ref{ch24:fig:review-magnetthreeloopsdiagram-1}中，轴向箭头描述对相对运动的阻碍；径向标记只适用于图示近轴极外和近中部的局部场条件。
\par\noindent\begin{minipage}{\linewidth}
\magnetThreeLoopsDiagram
\figcap{磁铁左移时的三个圆环}\label{ch24:fig:review-magnetthreeloopsdiagram-1}
\end{minipage}\par\medskip
\end{xesolution}

\Needspace{6\baselineskip}
\begin{origexamplebox}{例题 9.3【小棒加电阻】}
一根导体棒横跨在宽度为 $L$ 的光滑导轨上，导轨间有垂直纸面向里的匀强磁场 $B$。棒的质量为 $m$，左侧电阻为 $R$，
导轨和棒的电阻不计。初始时棒以速度 $v_0$ 向右运动。

(1) 求任意时刻导体棒的速度；

(2) 求从开始运动到 $t\to\infty$ 通过电阻的总电量。

% Pure source-question geometry: no derived forces, currents or solution curves.
%
\begin{center}\begin{tikzpicture}[scale=0.78,>=Latex]
  \draw[Ink,very thick] (-3,1.4)--(3,1.4); \draw[Ink,very thick] (-3,-1.4)--(3,-1.4);
  \draw[Ink,very thick] (-2.7,-1.4)--(-2.7,-0.45) decorate[decoration={zigzag,segment length=4pt,amplitude=2pt}] {-- (-2.7,0.45)}--(-2.7,1.4);
  \node[left=8pt] at (-2.7,0) {$R$};
  \draw[DeepBlue,very thick] (0.8,-1.55)--(0.8,1.55); \node[above] at (0.8,1.55) {$m$};
  \draw[vector] (1.0,0)--(2.4,0) node[right] {$v_0$};
  \draw[<->] (-3.7,-1.4)--node[left] {$L$}(-3.7,1.4);
  \foreach \x in {-1.7,-0.7,0.3,1.4,2.3}{\foreach \y in {-0.75,0.75}{\node[TextGray] at (\x,\y) {$\times$};}}
  \end{tikzpicture}\end{center}
\figcap{棒与定值电阻组成的闭合导轨回路。}\label{fig:em24-question-3}
\end{origexamplebox}

\begin{xesolution}
棒向右运动时回路面积增大，向里磁通增大。感应电流的 Ampere 力必然向左。当速度为 $v$ 时，
\begin{equation}
  \mathcal E=BLv,
  \qquad I=\frac{BL}{R}v,
  \qquad m\dot v=-BLI=-\frac{B^2L^2}{R}v.
\end{equation}
令
\begin{equation}
  \tau=\frac{mR}{B^2L^2},
\end{equation}
得
\begin{equation}
  \boxed{v(t)=v_0\e^{-t/\tau}}.
\end{equation}
总电量可以直接积分，也可用冲量消去时间：
\begin{equation}
  q_{\rm tot}=\int_0^\infty I\dd t
  =\frac{BL}{R}v_0\tau
  =\boxed{\frac{mv_0}{BL}}.
\end{equation}
结果与 $R$ 无关：电阻越大，瞬时电流越小，但衰减时间按同样比例变长。

图~\ref{ch24:fig:review-rodresistordiagram-1}中，棒横跨上下导轨，左侧R闭合回路；速度减小导致电流和阻力同时减小。
\par\noindent\begin{minipage}{\linewidth}
\rodResistorDiagram
\figcap{导体棒经电阻耗散动能}\label{ch24:fig:review-rodresistordiagram-1}
\end{minipage}\par\medskip
\end{xesolution}

\Needspace{6\baselineskip}
\begin{origexamplebox}{例题 9.4【小棒加小棒】}
两根导体棒横跨在宽度为 $L$ 的光滑导轨上，导轨间有垂直纸面向里的匀强磁场 $B$。两棒质量均为 $m$、电阻均为 $R$，
导轨电阻不计。右棒初速为 $v_0$ 向右，左棒初始静止。

(1) 求任意时刻左右两棒的速度；

(2) 求两棒从开始运动到 $t\to\infty$ 的相对位移。

% Pure source-question geometry: no derived forces, currents or solution curves.
%
\begin{center}\begin{tikzpicture}[scale=0.78,>=Latex]
  \draw[Ink,very thick] (-3,1.4)--(3,1.4); \draw[Ink,very thick] (-3,-1.4)--(3,-1.4);
  \draw[DeepBlue,very thick] (-1.5,-1.55)--(-1.5,1.55); \draw[DeepBlue,very thick] (1.5,-1.55)--(1.5,1.55);
  \node[above] at (-1.5,1.62) {$m,R$}; \node[above] at (1.5,1.62) {$m,R$};
  \draw[vector] (1.7,0)--(2.9,0) node[right] {$v_0$};
  \foreach \x in {-2.4,-0.5,0.5,2.4}{\foreach \y in {-0.7,0.7}{\node[TextGray] at (\x,\y) {$\times$};}}
  \draw[<->] (-3.25,-1.4)--node[left] {$L$}(-3.25,1.4);
  \end{tikzpicture}\end{center}
\figcap{两根相同导体棒，右棒具有初速度。}\label{fig:em24-question-4}
\end{origexamplebox}

\begin{xesolution}
记左、右棒速度为 $v_1,v_2$，相对速度 $u=v_2-v_1$。回路电动势只与面积变化率有关：
\begin{equation}
  \mathcal E=BLu,
  \qquad I=\frac{BL}{2R}u.
\end{equation}
左棒受力向右，右棒受力向左，故
\begin{equation}
  m\dot v_1=BLI,
  \qquad m\dot v_2=-BLI.
\end{equation}
因而质心速度不变，$V_{\rm cm}=v_0/2$，而
\begin{equation}
  \dot u=-\frac{B^2L^2}{mR}u.
\end{equation}
令 $\tau=mR/(B^2L^2)$，得
\begin{equation}
  \boxed{v_1=\frac{v_0}{2}\left(1-\e^{-t/\tau}\right)},
  \qquad
  \boxed{v_2=\frac{v_0}{2}\left(1+\e^{-t/\tau}\right)}.
\end{equation}
两棒的相对位移为
\begin{equation}
  \boxed{\Delta x_{\rm rel}=\int_0^\infty u\dd t
  =v_0\tau=\frac{mRv_0}{B^2L^2}}.
\end{equation}

图~\ref{ch24:fig:review-tworoddiagram-1}中，两个运动边界的电动势相减，驱动量取决于 $v_2-v_1$，不是任一棒的绝对速度。
\par\noindent\begin{minipage}{\linewidth}
\twoRodDiagram
\figcap{同一导轨上的两根导体棒}\label{ch24:fig:review-tworoddiagram-1}
\end{minipage}\par\medskip
\end{xesolution}

\Needspace{6\baselineskip}
\begin{origexamplebox}{例题 9.5【小棒再加小棒】}
两根导体棒分别横跨在左右两组导轨上，左侧导轨宽度为 $L_1$、磁场为 $B_1$，右侧导轨宽度为 $L_2$、磁场为 $B_2$，
中间用很长的过渡导轨连接。两棒质量分别为 $m_1,m_2$，回路总电阻为 $R$。右棒初速为 $v_0$ 向右，左棒初始静止。

(1) 求任意时刻两棒的速度；

(2) 找出运动中的一个守恒量。

% Pure source-question geometry: no derived forces, currents or solution curves.
%
\begin{center}\begin{tikzpicture}[scale=.82,>=Latex]
 \draw[Ink,thick](-5,1)--(-1,1)--(1,1.8)--(5,1.8);
 \draw[Ink,thick](-5,-1)--(-1,-1)--(1,-1.8)--(5,-1.8);
 \draw[DeepBlue,very thick](-4,-1.15)--(-4,1.15)node[above]{$m_1$};
 \draw[DeepBlue,very thick](4,-1.95)--(4,1.95)node[above]{$m_2$};
 \draw[vector](4.25,0)--(5.4,0)node[right]{$v_0$};
 \node at(-2,-.5){$B_1,L_1$};\node at(2,-.6){$B_2,L_2$};
 \foreach \x in {-3,-2,2,3}{\node[TextGray]at(\x,.65){$\times$};}
 \end{tikzpicture}\end{center}
\figcap{不同棒长与场强的两导轨段，由固定导线相连。}\label{fig:em24-question-5}
\end{origexamplebox}

\begin{xesolution}
记 $k_1=B_1L_1$、$k_2=B_2L_2$。按图中回路方向取正，瞬时感应电动势为
\begin{equation}
  \mathcal E=k_2v_2-k_1v_1,
  \qquad I=\frac{k_2v_2-k_1v_1}{R}.
\end{equation}
两棒的方程是
\begin{equation}
  m_1\dot v_1=k_1I,
  \qquad m_2\dot v_2=-k_2I.
\end{equation}
于是
\begin{equation}
  \frac{\dd}{\dd t}\left(\frac{m_1v_1}{k_1}+\frac{m_2v_2}{k_2}\right)=0,
\end{equation}
所求守恒量可取为
\begin{equation}
  \boxed{\frac{m_1v_1}{B_1L_1}+\frac{m_2v_2}{B_2L_2}
  =\frac{m_2v_0}{B_2L_2}}.
\end{equation}
再令 $y=k_2v_2-k_1v_1$，有
\begin{equation}
  \dot y=-\frac1R\left(\frac{k_1^2}{m_1}+\frac{k_2^2}{m_2}\right)y,
  \qquad
  y(0)=k_2v_0.
\end{equation}
定义
\begin{equation}
  D=k_1^2m_2+k_2^2m_1,
  \qquad \tau=\frac{Rm_1m_2}{D},
\end{equation}
积分得
\begin{align}
  \boxed{v_1(t)}&=\boxed{\frac{k_1k_2m_2}{D}v_0\left(1-\e^{-t/\tau}\right)},\\
  \boxed{v_2(t)}&=\boxed{\frac{k_1^2m_2+k_2^2m_1\e^{-t/\tau}}{D}v_0}.
\end{align}
当 $t\to\infty$ 时 $k_2v_2=k_1v_1$，也就是总动生电动势降为零。

图~\ref{ch24:fig:review-unequalroddiagram-1}中，上下两根过渡导线分别连通；右棒正电流向上、左棒向下，故两棒的安培力方向相反。
\par\noindent\begin{minipage}{\linewidth}
\unequalRodDiagram
\figcap{不等宽导轨之间的闭合连接}\label{ch24:fig:review-unequalroddiagram-1}
\end{minipage}\par\medskip
\end{xesolution}

\Needspace{6\baselineskip}
\begin{origexamplebox}{例题 9.6【小棒加电容】}
一根质量为 $m$ 的导体棒横跨在宽度为 $L$ 的光滑导轨上，导轨间有垂直纸面向里的匀强磁场 $B$。回路中串联电容 $C$，
并保留串联小电阻 $R_s>0$；忽略回路自感及其他电阻。在 $t=0$ 接通回路，棒初速为 $v_0$ 向右，电容初始不带电。
求充电暂态中的棒速和电荷、长时间后的终态，并讨论先等待暂态结束再取 $R_s\to0^+$ 的极限。

% Pure source-question geometry: no derived forces, currents or solution curves.
%
\begin{center}\begin{tikzpicture}[scale=0.78,>=Latex]
  \draw[Ink,very thick] (-3,1.4)--(-1.9,1.4);
  \draw[Ink,thick] (-1.9,1.2) rectangle (-1.0,1.6);
  \node[above] at (-1.45,1.6) {$R_s$};
  \draw[Ink,very thick] (-1.0,1.4)--(3,1.4); \draw[Ink,very thick] (-3,-1.4)--(3,-1.4);
  \draw[Ink,very thick] (-2.7,-1.4)--(-2.7,-0.25); \draw[Ink,very thick] (-3.05,-0.25)--(-2.35,-0.25);
  \draw[Ink,very thick] (-3.05,0.25)--(-2.35,0.25); \draw[Ink,very thick] (-2.7,0.25)--(-2.7,1.4);
  \node[left] at (-2.7,0) {$C$};
  \draw[DeepBlue,very thick] (0.8,-1.55)--(0.8,1.55) node[above] {$m$};
  \draw[vector] (1.0,0)--(2.4,0) node[right] {$v_0$};
  \foreach \x in {-1.7,-0.7,0.2,1.5,2.4}{\foreach \y in {-0.7,0.7}{\node[TextGray] at (\x,\y) {$\times$};}}
  \end{tikzpicture}\end{center}
\figcap{导体棒通过小电阻与初始未充电的电容器相连，图示接通后的回路。}\label{fig:em24-question-6}
\end{origexamplebox}

\begin{xesolution}
接通回路后先经历充电暂态。先求长时间极限，再解出完整的时间变化。
设长时间后电容带电量为 $q$、棒速为 $v$。此时电流趋于零，回路方程与暂态全过程的冲量方程分别为
\begin{equation}
  \frac qC=BLv,
  \qquad m(v-v_0)=-BLq.
\end{equation}
解得
\begin{equation}
  \boxed{v=\frac{m}{m+CB^2L^2}v_0},
  \qquad
  \boxed{q=\frac{mCBL}{m+CB^2L^2}v_0}.
\end{equation}
上述 $q,v$ 分别记作 $q_\infty,v_\infty$；这是题设 $R_s>0$ 时的长时间极限，电流渐近趋零，棒速渐近趋于常量。
为明确过程，令 $k_m=BL$，则
\[
 R_s\dot q+q/C=k_mv,\qquad mv+k_mq=mv_0.
\]
消去 $v$ 得 $R_s\dot q+(1/C+k_m^2/m)q=k_mv_0$。若 $q(0)=0$，则
\[
 q(t)=q_\infty(1-\e^{-t/\tau_s}),\quad
 \tau_s=\frac{R_smC}{m+Ck_m^2},\quad
 v(t)=v_0-\frac{k_m}{m}q(t).
\]
因此先令 $t\gg\tau_s$ 再取 $R_s\to0^+$，终态不变，暂态持续时间趋零。
能量检验给出
\[
 Q=\frac12mv_0^2-\frac12mv_\infty^2-\frac{q_\infty^2}{2C}
 =\frac{mCk_m^2v_0^2}{2(m+Ck_m^2)}>0.
\]
这个耗散不随上述顺序下的 $R_s\to0^+$ 而消失；直接令电阻、自感都严格为零会得到奇异初值问题。
若保留电感却使电阻严格为零，系统一般振荡，不能说它必然进入同一恒速终态。

图~\ref{ch24:fig:review-rodcapacitordiagram-1}中，电容最终电压等于棒的动生电动势；暂态需按文中微小电阻模型理解。
\par\noindent\begin{minipage}{\linewidth}
\rodCapacitorDiagram
\figcap{电容充电时的回路与棒的冲量关系}\label{ch24:fig:review-rodcapacitordiagram-1}
\end{minipage}\par\medskip
\end{xesolution}

\Needspace{6\baselineskip}
\begin{origexamplebox}{例题 9.7【小棒加电感】}
一根质量为 $m$ 的导体棒横跨在宽度为 $\ell$ 的光滑导轨上，导轨间有垂直纸面向里的匀强磁场 $B$。回路中串联自感为 $\mathcal L$
的理想电感，忽略回路其余部分的自感与电阻，$\mathcal L$ 不随棒的位置变化。棒初速为 $v_0$ 向右，电感初始无电流。描述 $t=0$ 后导体棒的运动。

% Pure source-question geometry: no derived forces, currents or solution curves.
%
\begin{center}\begin{tikzpicture}[scale=0.78,>=Latex]
  \draw[Ink,very thick] (-3,1.4)--(3,1.4); \draw[Ink,very thick] (-3,-1.4)--(3,-1.4);
  \draw[Ink,very thick] (-2.7,-1.4)--(-2.7,-0.6);
  \draw[Ink,thick] (-2.7,-0.6) decorate[decoration={coil,aspect=0.65,segment length=5pt,amplitude=5pt}] {-- (-2.7,0.6)};
  \draw[Ink,very thick] (-2.7,0.6)--(-2.7,1.4); \node[left] at (-2.9,0) {$\mathcal L$};
  \draw[DeepBlue,very thick] (0.8,-1.55)--(0.8,1.55) node[above] {$m$};
  \draw[vector] (1.0,0)--(2.4,0) node[right] {$v_0$};
  \foreach \x in {-1.7,-0.7,0.2,1.5,2.4}{\foreach \y in {-0.7,0.7}{\node[TextGray] at (\x,\y) {$\times$};}}
  \end{tikzpicture}\end{center}
\figcap{导体棒与电感相连，电感初始电流为零。}\label{fig:em24-question-7}
\end{origexamplebox}

\begin{xesolution}
取感应电流 $I$ 的正方向使它对棒产生向左的力。电路与力学方程为
\begin{equation}
  \mathcal L\dot I=B\ell v,
  \qquad m\dot v=-B\ell I.
\end{equation}
对第二式再求导并代入第一式，得
\begin{equation}
  \ddot v+\Omega^2v=0,
  \qquad
  \Omega=\frac{B\ell}{\sqrt{m\mathcal L}}.
\end{equation}
用 $v(0)=v_0$、$I(0)=0$ 得
\begin{equation}
  \boxed{v(t)=v_0\cos\Omega t},
  \qquad
  \boxed{I(t)=v_0\sqrt{\frac m{\mathcal L}}\sin\Omega t},
\end{equation}
\begin{equation}
  \boxed{x(t)-x(0)=\frac{v_0}{\Omega}\sin\Omega t}.
\end{equation}
棒的动能和电感磁能周期性交换：
\begin{equation}
  \frac12mv^2+\frac12\mathcal LI^2=\frac12mv_0^2.
\end{equation}

图~\ref{ch24:fig:review-rodinductordiagram-1}中，线圈接入回路后，电感电流和棒速度连续变化，动能与磁能相互转化。
\par\noindent\begin{minipage}{\linewidth}
\rodInductorDiagram
\figcap{导体棒与电感的能量交换}\label{ch24:fig:review-rodinductordiagram-1}
\end{minipage}\par\medskip
\end{xesolution}

\Needspace{6\baselineskip}
\begin{origexamplebox}{例题 9.8【小棒加电源】}
一根质量为 $m$ 的导体棒横跨在宽度为 $\ell$ 的光滑导轨上，导轨间有垂直纸面向里的匀强磁场 $B$。回路总电阻为 $R$，
并接入恒定电源 $\mathcal E$，电源极性使静止棒受到向右的安培力。棒初速为 $v_0$ 向右。描述 $t=0$ 后导体棒的运动。

% Pure source-question geometry: no derived forces, currents or solution curves.
%
\begin{center}\begin{tikzpicture}[scale=.8,>=Latex]
 \draw[Ink,thick](-3,1.5)--(3,1.5) (-3,-1.5)--(3,-1.5);
 \draw[Ink,thick](-2.7,-1.5)--(-2.7,-.8);
 \draw[Ink,thick](-2.9,-.8)rectangle(-2.5,-.15);
 \draw[Ink,thick](-2.7,-.15)--(-2.7,.25);
 \draw[Ink,thick](-2.95,.25)--(-2.45,.25) (-3.1,.65)--(-2.3,.65);
 \draw[Ink,thick](-2.7,.65)--(-2.7,1.5);
 \node[left]at(-3,-.5){$R$};\node[left]at(-3.15,.45){$\mathcal E$};
 \node[right]at(-2.35,.25){$-$};\node[right]at(-2.25,.65){$+$};
 \draw[DeepBlue,very thick](.8,-1.65)--(.8,1.65)node[above]{$m$};
 \draw[vector](1.2,0)--(2.6,0)node[right]{$v_0$};
 \foreach \x in {-1.6,-.6,1.8}{\foreach \y in {-.9,.9}{\node[TextGray]at(\x,\y){$\times$};}}
\end{tikzpicture}\end{center}
\figcap{电源、定值电阻和运动导体棒的实际连接。}\label{fig:em24-question-8}
\end{origexamplebox}

\begin{xesolution}
取棒中向下为电流正方向。按右手定则，正电流在向纸内的磁场中产生向右的安培力；
实际电流的正负还取决于棒速，不能仅由电源极性判断。按此约定，电流和运动方程为
\begin{equation}
  I=\frac{\mathcal E-B\ell v}{R},
  \qquad
  m\dot v=B\ell I.
\end{equation}
运动电动势 $B\ell v$ 是电源的反电动势。定义
\begin{equation}
  v_\infty=\frac{\mathcal E}{B\ell},
  \qquad \tau=\frac{mR}{B^2\ell^2},
\end{equation}
得
\begin{equation}
  \boxed{v(t)=v_\infty+(v_0-v_\infty)\e^{-t/\tau}}.
\end{equation}
若 $v_0<v_\infty$，电源使棒加速；若 $v_0>v_\infty$，反电动势大于电源，电流反向，棒减速；两种情况都趋向同一终端速度。

图~\ref{ch24:fig:review-rodsourcediagram-1}中，按所标电源极性，正电流在棒中向下；当动生电动势大于电源时，电流可反向。
\par\noindent\begin{minipage}{\linewidth}
\rodSourceDiagram
\figcap{电源驱动的导体棒}\label{ch24:fig:review-rodsourcediagram-1}
\end{minipage}\par\medskip
\end{xesolution}

\Needspace{6\baselineskip}
\begin{origexamplebox}{例题 9.9【小棒大结局】}
一根质量为 $m$、电阻为 $R$ 的导体棒横跨在宽度为 $L$ 的光滑导轨上，导轨间有垂直纸面向里的匀强磁场 $B$。

(1) 若 $R=0$，推导棒两端电压 $U$ 与流过棒的电流 $I$ 的关系，辨认其等效元件与参数；

(2) 若 $R\ne0$，再求 $U$--$I$ 关系与等效元件；

(3) 用该结论重新理解前面的电阻、电容、电感和电源问题。

% Pure source-question geometry: no derived forces, currents or solution curves.
\begin{center}\begin{tikzpicture}[scale=.78,>=Latex]
 \draw[thick](-3,1.4)--(3,1.4) (-3,-1.4)--(3,-1.4);
 \draw[DeepBlue,very thick](.8,-1.55)--(.8,1.55)node[above]{$m,R$};
 \draw[vector](1.1,0)--(2.4,0)node[right]{$v$};
 \draw[vector](-2.7,1.4)--(-1.7,1.4)node[above]{$I$};
 \node[left]at(-3,1.4){$+$};\node[left]at(-3,-1.4){$-$};
 \draw[<->](-3.7,-1.4)--(-3.7,1.4)node[midway,left]{$U$};
 \draw[<->](3.4,-1.4)--(3.4,1.4)node[midway,right]{$L$};
 \foreach \x in {-1.6,-.5,1.8}{\foreach \y in {-.8,.8}{\node at(\x,\y){$\otimes$};}}
 \node at(-.7,0){$B$};
\end{tikzpicture}\end{center}
\figcap{可运动导体棒的电学端口；箭头和正负号定义关联参考方向。}\label{fig:em24-question-9}
\end{origexamplebox}

\begin{xesolution}
把棒的动生电动势记为 $U_m=BLv$，取正电流的 Ampere 力使棒向右，则
\begin{equation}
  m\dot v=BLI.
\end{equation}
因而
\begin{equation}
  \dot U_m=BL\dot v=\frac{B^2L^2}{m}I,
  \qquad
  \boxed{I=C_m\dot U_m},
  \qquad
  \boxed{C_m=\frac{m}{B^2L^2}}.
\end{equation}
当 $R=0$ 时，$U=U_m$，所以这根可运动导体棒在电学端口上等效为电容 $C_m$。其“电场能”正好就是棒的动能：
\begin{equation}
  \frac12C_mU_m^2=\frac12\frac{m}{B^2L^2}(BLv)^2=\frac12mv^2.
\end{equation}
当棒有电阻 $R$ 时，端口方程为
\begin{equation}
  \boxed{U=RI+U_m},
  \qquad I=C_m\dot U_m,
\end{equation}
因而等效为电阻 $R$ 与电容 $C_m$ 串联，电容初电压为 $U_m(0)=BLv_0$。

前面诸问题因此统一为基本电路：例 9.3 是 $RC_m$ 放电；例 9.6 是机械等效电容与真实电容的瞬态电荷重分配；
例 9.7 是 $\mathcal LC_m$ 振荡；例 9.8 是电源对 $RC_m$ 的充电。

图~\ref{ch24:fig:review-equivalentcircuitdiagram-1}中，机械端口电压 $U_m=BLv$ 与电流关系相当于电容，图中R与该电容串联。
\par\noindent\begin{minipage}{\linewidth}
\equivalentCircuitDiagram
\figcap{运动导体棒的电学端口等效}\label{ch24:fig:review-equivalentcircuitdiagram-1}
\end{minipage}\par\medskip
\end{xesolution}

\Needspace{6\baselineskip}
\begin{origexamplebox}{例题 9.10【小环的响应】}
半径为 $R$ 的固定绝缘大圆环与匀强磁场垂直，一个质量为 $m$、电荷为 $+q$ 的小环套在大环上，可无摩擦滑动。规定绕“垂直纸面
向里”的轴按右手螺旋法则为角坐标 $\theta$ 正向，初始 $\theta=0$ 且静止。

(1) $B(t)=\alpha t$；

(2) $B(t)=B_0\sin\omega t$；

(3) $B(t)=B_0\cos\omega t$。

分别求小环的角速度、位置以及大环对小环的径向支持力。
\par\noindent\begin{minipage}{\linewidth}
\chargedRingDiagram
\figcap{小环的角向和径向符号约定}\label{ch24:fig:review-chargedringdiagram-1}
\end{minipage}\par\medskip
\end{origexamplebox}

\begin{xesolution}
半径 $R$ 处的感生电场为
\begin{equation}
  E_\theta=-\frac R2\dot B.
\end{equation}
切向方程 $mR\ddot\theta=qE_\theta$ 给出一个对任意 $B(t)$ 都成立的通式：
\begin{equation}
  \boxed{\dot\theta(t)=-\frac q{2m}\,[B(t)-B(0)]},
  \qquad
  \boxed{\theta(t)=-\frac q{2m}\int_0^t[B(t')-B(0)]\dd t'}.
\end{equation}
取径向向外为正。磁力的径向分量为 $qR\dot\theta B$，径向加速度为 $-R\dot\theta^2$，故大环的有符号支持力为
\begin{equation}
  \boxed{N_r=-mR\dot\theta^2-qR\dot\theta B}.
\end{equation}
$N_r>0$ 表示支持力向外，$N_r<0$ 表示向内。

三种情况依次为
\begin{align}
  B=\alpha t:\quad
  &\boxed{\dot\theta=-\frac{q\alpha}{2m}t},
  &\boxed{\theta=-\frac{q\alpha}{4m}t^2};\\
  B=B_0\sin\omega t:\quad
  &\boxed{\dot\theta=-\frac{qB_0}{2m}\sin\omega t},
  &\boxed{\theta=\frac{qB_0}{2m\omega}(\cos\omega t-1)};\\
  B=B_0\cos\omega t:\quad
  &\boxed{\dot\theta=\frac{qB_0}{2m}(1-\cos\omega t)},
  &\boxed{\theta=\frac{qB_0}{2m}\left(t-\frac{\sin\omega t}{\omega}\right)}.
\end{align}
对第 (1) 问，将 $\dot\theta=-q\alpha t/(2m)$ 代入，向心项为
$-q^2R\alpha^2t^2/(4m)$，磁力移项为 $q^2R\alpha^2t^2/(2m)$，相加得
\begin{equation}
  \boxed{N_r=\frac{q^2R\alpha^2t^2}{4m}},
\end{equation}
方向向外：此时磁力向内且大于所需向心力，大环必须向外拉住小环。
对第 (2) 问同理得
\[
 N_r=-\frac{q^2RB_0^2}{4m}\sin^2\omega t
      +\frac{q^2RB_0^2}{2m}\sin^2\omega t
 =\boxed{\frac{q^2RB_0^2}{4m}\sin^2\omega t}.
\]
它向外或为零。对第 (3) 问，令 $c=\cos\omega t$，得到
\[
 N_r=-\frac{q^2RB_0^2}{4m}[(1-c)^2+2c(1-c)]
 =\boxed{-\frac{q^2RB_0^2}{4m}\sin^2\omega t}.
\]
它向内或为零。两者不同的原因是初始磁场分别为0与 $B_0$，角动量的积分常数不同，
不能只凭某时刻 $B$ 的大小判断速度。三问均满足初始速度为零；$N_r$ 的量纲为力。

图~\ref{ch24:fig:review-chargedringdiagram-1}中，顺时针为角坐标正向，径向向外为正；两个力的代数值按同一径向约定相加。

\end{xesolution}
